% ============================================================
%  ON THE PARTITION ALGEBRA OF BOUNDED PHASE SPACE:
%  COMPOSITION-INFLATION, ANGULAR RESOLUTION, AND THE
%  CATEGORICAL REFORMULATION OF FUNDAMENTAL CONSTANTS
% ============================================================

\documentclass[twocolumn,10pt,a4paper]{article}

\usepackage[utf8]{inputenc}
\usepackage[T1]{fontenc}
\usepackage{lmodern}
\usepackage[a4paper,top=2cm,bottom=2.2cm,left=1.8cm,right=1.8cm,columnsep=0.6cm]{geometry}
\usepackage{amsmath,amssymb,amsthm,mathtools}
\usepackage{bm}
\usepackage{booktabs}
\usepackage{array}
\usepackage{enumitem}
\usepackage{graphicx}
\usepackage{xcolor}
\usepackage{microtype}
\usepackage[round,sort&compress]{natbib}
\usepackage[colorlinks=true,linkcolor=blue!60!black,
            citecolor=blue!60!black,urlcolor=blue!60!black]{hyperref}
\usepackage{fancyhdr}

\pagestyle{fancy}
\fancyhf{}
\fancyhead[LE,RO]{\thepage}
\fancyhead[RE]{\small Partition Algebra of Bounded Phase Space}
\fancyhead[LO]{\small Sachikonye}
\renewcommand{\headrulewidth}{0.4pt}

% ---- Theorem environments ----------------------------------------
\newtheoremstyle{thmstyle}{\topsep}{\topsep}{\itshape}{}{\bfseries}{.}{.5em}{}
\theoremstyle{thmstyle}
\newtheorem{theorem}{Theorem}[section]
\newtheorem{lemma}[theorem]{Lemma}
\newtheorem{proposition}[theorem]{Proposition}
\newtheorem{corollary}[theorem]{Corollary}
\theoremstyle{definition}
\newtheorem{definition}[theorem]{Definition}
\newtheorem{axiom}{Axiom}
\theoremstyle{remark}
\newtheorem{remark}[theorem]{Remark}

% ---- Custom commands --------------------------------------------
\newcommand{\R}{\mathbb{R}}
\newcommand{\N}{\mathbb{N}}
\newcommand{\Zp}{\mathbb{Z}^{+}}
\newcommand{\tP}{t_{\mathrm{P}}}
\newcommand{\lP}{\ell_{\mathrm{P}}}
\newcommand{\kB}{k_{\mathrm{B}}}
\newcommand{\Parts}{\mathcal{P}}

% =========================================================
\title{\textbf{On the Partition Algebra of Bounded Phase Space:\\
Composition-Inflation, Angular Resolution, and\\
the Categorical Reformulation of Fundamental Constants}}

\author{
Kundai Farai Sachikonye\\
Technical University of Munich\\
Chair of Proteomics and Bioanalytics\\
\texttt{kundai.sachikonye@wzw.tum.de}
}
\date{\today}
% =========================================================

\begin{document}
\maketitle

% =========================================================
\begin{abstract}
\noindent
Every persistent physical system occupies a bounded region of phase space.
From this single empirical observation, promoted to an axiom, we derive by
deduction the complete partition algebra governing distinguishable states in
bounded oscillatory systems.

The central result is the \emph{composition-inflation formula}:
a bounded oscillatory system with $d$ entropy dimensions completing $n$
oscillation cycles generates
\begin{equation*}
T(n,d) = d \cdot (d+1)^{n-1}
\end{equation*}
categorically distinguishable trajectories — an exponential, not linear,
count. The formula is exact, has no adjustable parameters, and follows
from the binomial theorem applied to labeled integer compositions.

We derive three consequences that are not available in any existing framework.
First, the \emph{angular resolution} $\Delta\theta = 2\pi/T(n,d)$ is
dimensionless and therefore not subject to the Planck time constraint, which
bounds temporal intervals measured in seconds but does not bound dimensionless
ratios. For any $\varepsilon > 0$ there exists a finite $n$ achieving
$\Delta\theta < \varepsilon$.
Second, all fundamental constants except the gravitational constant $G$
reduce to tick counts, factors of $2\pi$, and a single energy unit
$E_{\mathrm{tick}}$ in categorical angular units: $c = 2\pi\,\mathrm{rad/tick}$,
$\hbar = E_{\mathrm{tick}}/(2\pi)$, $k_B = E_{\mathrm{tick}}/\ln(d+1)$.
Third, $G$ is the unique irreducible coupling constant — the only physical
constant not expressible in terms of counting and geometry.

The paper is self-contained: all results are proved from the single axiom
with no additional physical postulates. Applications to the LHC trigger,
mass spectrometry, and gravitational physics are indicated in the final section
but not developed here.

\medskip
\noindent\textbf{Keywords:} bounded phase space; partition algebra;
composition-inflation; angular resolution; Planck depth;
categorical angular units; gravitational constant
\end{abstract}

\tableofcontents

% =========================================================
\section{Introduction}
\label{sec:introduction}
% =========================================================

\subsection{The single empirical premise}

Physical systems exist. Every physical system that persists — an electron
bound to a nucleus, a molecule in a container, a planet in a stellar orbit,
an ion in a trap — occupies a bounded region of phase space. The complement
— a system whose phase space region grows without bound — disperses and ceases
to persist. This is not a theoretical proposal; it is a summary of every
observation of every persistent object in the history of physics.

We promote this observation to an axiom and derive consequences.

\subsection{What follows}

The derivation proceeds in four steps. Section~\ref{sec:foundations} states
the axiom and derives oscillatory necessity (Poincar\'e recurrence eliminates
every alternative). Section~\ref{sec:partition} constructs the partition
coordinate system $(n,\ell,m,s)$ with capacity $C(n) = 2n^2$ from
the geometry of bounded oscillation. Section~\ref{sec:inflation} derives
the composition-inflation formula $T(n,d) = d(d+1)^{n-1}$ and establishes
the angular resolution theorem. Section~\ref{sec:constants} expresses all
fundamental constants in categorical angular units, identifies $G$ as the
unique irreducible constant, and derives the Planck depth.

No theorem in Sections~\ref{sec:partition}--\ref{sec:constants} uses any
postulate beyond the axiom of Section~\ref{sec:foundations}.

% =========================================================
\section{Bounded Phase Space and Oscillatory Necessity}
\label{sec:foundations}
% =========================================================

\begin{axiom}[Bounded Phase Space]
\label{ax:bounded}
Every persistent dynamical system $(\Omega, \mu, \phi_t)$ satisfies:
\begin{enumerate}[nosep,label=(\roman*)]
\item $\Omega \subset \R^{2d}$ is bounded with finite Liouville measure
$\mu(\Omega) < \infty$.
\item $\phi_t : \Omega \to \Omega$ preserves $\mu$ for all $t$.
\item $\Omega$ admits a sequence of finite measurable partitions
$\mathcal{P}_1 \preceq \mathcal{P}_2 \preceq \cdots$ with each $\mathcal{P}_{k+1}$
refining $\mathcal{P}_k$ and the generated $\sigma$-algebra equalling the
Borel algebra modulo null sets.
\end{enumerate}
\end{axiom}

\begin{theorem}[Poincar\'e Recurrence \citealp{poincare1890}]
\label{thm:poincare}
Under Axiom~\ref{ax:bounded}, for every measurable $A \subset \Omega$ with
$\mu(A) > 0$ and every $T > 0$, almost every $x \in A$ satisfies
$\phi_{t_n}(x) \in A$ for infinitely many $t_n \to \infty$.
\end{theorem}

\begin{proof}
Suppose not. Let $B = \{x \in A : \phi_t(x) \notin A\;\forall\, t > T\}$.
The iterates $\{\phi_{-kT}(B)\}_{k \geq 0}$ are pairwise disjoint (if
$x \in \phi_{-jT}(B) \cap \phi_{-kT}(B)$ with $j < k$, then
$\phi_{(k-j)T}(x) \in B$ and later $\phi_t(x)$ returns to $A$,
contradicting $x \in B$). Measure preservation gives
$\sum_{k=0}^\infty \mu(B) = \mu(\bigcup_k \phi_{-kT}(B)) \leq \mu(\Omega) < \infty$,
so $\mu(B) = 0$.
\end{proof}

\begin{corollary}[Oscillatory Necessity]
\label{cor:oscillatory}
Under Axiom~\ref{ax:bounded}, every persistent non-static trajectory is
oscillatory: no coordinate can grow monotonically, and every trajectory
returns to every neighbourhood of its starting point infinitely often.
\end{corollary}

\begin{proof}
Monotonic growth contradicts boundedness. Absence of recurrence contradicts
Theorem~\ref{thm:poincare}. The only consistent dynamics is recurrent
oscillatory motion.
\end{proof}

For integrable systems, the Arnold--Liouville theorem \citep{arnold1989}
guarantees action-angle variables $(J_i, \theta_i)$ with quasi-periodic
motion $\dot\theta_i = \omega_i(J)$. The Koopman operator \citep{koopman1931}
on $L^2(\Omega,\mu)$ has a pure point spectrum on cyclic subspaces; its
eigenfrequencies $\{\omega_k\}$ are the mode frequencies of the system.
Near-integrable systems retain this structure on most tori by the
KAM theorem \citep{kolmogorov1954,arnold1963,moser1962}.

% =========================================================
\section{Partition Coordinates and Shell Capacity}
\label{sec:partition}
% =========================================================

\subsection{Hierarchical partition and coordinate generation}

The hierarchical partition of Axiom~\ref{ax:bounded}(iii) is not merely
formal. For a three-dimensional bounded system with $SO(3)$ symmetry, the
refinement sequence generates a natural four-coordinate label on each
partition cell.

\begin{definition}[Partition State]
A \emph{partition state} of a bounded three-dimensional Hamiltonian system
is an equivalence class of phase-space trajectories sharing the same closure
label $(n, \ell, m, s)$ where:
$n \in \{1,2,3,\ldots\}$ (principal level, radial closure count),
$\ell \in \{0,1,\ldots,n-1\}$ (angular sublevel, polar closure count),
$m \in \{-\ell,\ldots,+\ell\}$ (azimuthal projection),
$s \in \{-\tfrac12, +\tfrac12\}$ (rotational parity).
The set of all partition states is $\Parts$.
\end{definition}

The label $(n,\ell,m,s)$ is derived from the closure conditions of orbits in
action-angle coordinates. An orbit closes when $n_r \omega_r = n_\theta
\omega_\theta = n_\phi \omega_\phi$ for positive integers; the hierarchy
of radial, polar, and azimuthal boundaries produces the nesting
$0 \leq \ell \leq n-1$, $|m| \leq \ell$, and the binary parity $s$
from the double cover $\mathrm{Spin}(3) \cong SU(2) \to SO(3)$
\citep{wigner1959}.

\begin{theorem}[Shell Capacity]
\label{thm:capacity}
The number of distinct partition states at principal level $n$ is
\begin{equation}
C(n) = 2n^2.
\label{eq:capacity}
\end{equation}
\end{theorem}

\begin{proof}
At level $n$, $\ell$ ranges over $\{0,\ldots,n-1\}$; for each $\ell$,
$m$ ranges over $2\ell+1$ values; $s$ contributes a factor of 2:
\begin{align}
C(n) = 2\sum_{\ell=0}^{n-1}(2\ell+1) = 2\bigl(n + 2\cdot\tfrac{n(n-1)}{2}\bigr)
= 2n^2.\qed
\end{align}
\end{proof}

\begin{corollary}[Cumulative Count]
The total number of partition states at or below level $n_{\max}$ is
\begin{equation}
N_{\mathrm{state}}(n_{\max}) = \sum_{n=1}^{n_{\max}} 2n^2
= \frac{n_{\max}(n_{\max}+1)(2n_{\max}+1)}{3}.
\end{equation}
\end{corollary}

\begin{remark}
The sequence $C(n) = 2, 8, 18, 32, 50, 72, 98, \ldots$ is the electron
shell capacity sequence of the periodic table. This is not coincidental:
both arise from the same $SO(3)$ representation theory applied to a bounded
three-dimensional oscillatory system.
\end{remark}

\subsection{The partition bijection}

\begin{theorem}[Partition Bijection]
\label{thm:bijection}
There exists a bijection $\Phi : \Zp \to \Parts$ mapping each positive integer
$M$ to a unique partition state $(n,\ell,m,s)$. The map is constructive:
given $M$, the principal level is $n = \lfloor\sqrt{M/2}\rfloor + 1$
(adjusted to the correct shell), $\ell$ and $m$ follow from the lexicographic
ordering within the shell, and $s = \pm\tfrac12$ from the parity index.
\end{theorem}

\begin{proof}
Injectivity: distinct $M$ values yield distinct shell positions in the
lexicographic ordering, hence distinct $(n,\ell,m,s)$ tuples.
Surjectivity: every valid tuple $(n,\ell,m,s)$ has a unique cumulative index
$M = N_{\mathrm{state}}(n-1) + \delta$ where $\delta$ is the within-shell
offset; this $M$ maps back to $(n,\ell,m,s)$ by construction.
\end{proof}

\begin{theorem}[Strict Monotonicity of the Partition Count]
\label{thm:monotone}
Let $M(t)$ be the partition count of a bounded oscillatory system with
oscillation frequency $f > 0$. Then $M(t) = ft$ exactly, and
$M(t_2) > M(t_1)$ whenever $t_2 > t_1$ and at least one cycle has completed.
\end{theorem}

\begin{proof}
Each completed oscillation cycle increments $M$ by exactly 1. The rate is
$dM/dt = f = \omega/(2\pi)$, exact by definition of the oscillation count.
Strict monotonicity follows since $f > 0$.
\end{proof}

% =========================================================
\section{Composition-Inflation and Angular Resolution}
\label{sec:inflation}
% =========================================================

\subsection{Integer compositions}

\begin{definition}[Composition]
An \emph{integer composition} of $n$ is an ordered sequence $(c_1,\ldots,c_k)$
of positive integers summing to $n$. Two compositions differing in order
are distinct.
\end{definition}

\begin{theorem}[Composition Count \citealp{andrews1976}]
\label{thm:compcount}
The total number of compositions of $n$ is $2^{n-1}$.
\end{theorem}

\begin{proof}
Place $n$ objects in a row; there are $n-1$ gaps. Each gap either contains
a divider or does not. The $2^{n-1}$ binary choices on the gaps are in
bijection with compositions. \qed
\end{proof}

The number of compositions into exactly $k$ parts is $\binom{n-1}{k-1}$.

\subsection{Dimensional labeling}

The S-entropy space is the unit cube $[0,1]^d$ with $d$ entropy dimensions.
For $d = 3$ (knowledge $S_k$, temporal $S_t$, evolution $S_e$), any
oscillatory trajectory through this space can be characterized by which
dimension it refines at each step.

\begin{definition}[Labeled Composition]
A \emph{labeled composition} of $n$ in $d$ dimensions is a pair
$(\mathbf{C}, \mathbf{L})$ where $\mathbf{C} = (c_1,\ldots,c_k)$ is a
composition of $n$ and $\mathbf{L} = (\ell_1,\ldots,\ell_k)$ with
$\ell_i \in \{1,\ldots,d\}$ specifies the entropy dimension refined during
the $i$-th part of the trajectory.
\end{definition}

\begin{theorem}[Composition-Inflation Formula]
\label{thm:inflation}
The total number of labeled compositions of $n$ in $d$ dimensions is
\begin{equation}
\boxed{T(n,d) = d \cdot (d+1)^{n-1}.}
\label{eq:inflation}
\end{equation}
\end{theorem}

\begin{proof}
Summing over all numbers of parts:
\begin{align}
T(n,d) &= \sum_{k=1}^{n} \binom{n-1}{k-1} d^k
= d \sum_{j=0}^{n-1} \binom{n-1}{j} d^j
= d(1+d)^{n-1},
\end{align}
where the last step is the binomial theorem with $x = d$, $m = n-1$.
\end{proof}

\begin{proposition}[Growth Properties]
\label{prop:growth}
$T(1,d) = d$;\;
$T(n+1,d) = (d+1)\cdot T(n,d)$ (each additional cycle multiplies by $d+1$);\;
$\log T(n,d) = \log d + (n-1)\log(d+1)$;\;
$T(n,d)/n \to \infty$ as $n \to \infty$.
\end{proposition}

\begin{corollary}[S-Entropy Space]
For $d = 3$: $T(n,3) = 3 \cdot 4^{n-1}$.
After $n = 56$ cycles, $T(56,3) = 3 \cdot 4^{55} \approx 3.8 \times 10^{33}$.
\end{corollary}

\subsection{Trajectory distinguishability}

\begin{theorem}[Position-Trajectory Duality]
\label{thm:duality}
In the ternary encoding of the unit cube $[0,1]^3$, the trit string encodes
both the final position and the trajectory taken to reach it. Two distinct
labeled compositions produce distinct trit strings and are therefore
categorically distinguishable.
\end{theorem}

\begin{proof}
Case 1: distinct underlying compositions. Different sequences of segment
lengths determine transitions at different depths, yielding distinct trit
strings.
Case 2: same composition, distinct labels. There exists index $i$ with
$\ell_i \neq \ell'_i$; the $i$-th segment refines different axes, producing
distinct trit strings by orthogonality of the dimensions.
Both cases yield distinct categorical states. \qed
\end{proof}

\begin{theorem}[Non-Demolition Distinguishability]
\label{thm:nondemolition}
Let $\hat{O}_{\mathrm{cat}}$ measure categorical trajectory labels and
$\hat{O}_{\mathrm{phys}}$ measure any physical observable (position, momentum,
energy). Then $[\hat{O}_{\mathrm{cat}}, \hat{O}_{\mathrm{phys}}] = 0$.
Measuring which trajectory was taken does not disturb the physical state.
\end{theorem}

\begin{proof}
$\hat{O}_{\mathrm{cat}}$ acts on the trajectory label — a topological property
of the path. $\hat{O}_{\mathrm{phys}}$ acts on the instantaneous phase-space
position. These act on complementary degrees of freedom: the trajectory label
is constant along the trajectory (a conserved quantity), hence commutes with
the Hamiltonian and with all physical observables it generates.
Since $\hat{O}_{\mathrm{cat}}$ is diagonal in the trajectory basis and
$\hat{O}_{\mathrm{phys}}$ acts on the physical space factor, their tensor
product structure guarantees $[\hat{O}_{\mathrm{cat}}, \hat{O}_{\mathrm{phys}}] = 0$.
\qed
\end{proof}

\subsection{Angular resolution and the Planck constraint}

\begin{definition}[Angular Resolution]
For $T(n,d)$ distinguishable trajectories distributed uniformly over
$[0,2\pi)$, the angular resolution is
\[
\Delta\theta(n,d) = \frac{2\pi}{T(n,d)} = \frac{2\pi}{d(d+1)^{n-1}}.
\]
\end{definition}

\begin{theorem}[No Planck Limit on Angular Resolution]
\label{thm:noplanck}
The Planck time $\tP = \sqrt{\hbar G/c^5} \approx 5.39 \times 10^{-44}\,\mathrm{s}$
bounds temporal intervals measured in seconds. The angular resolution
$\Delta\theta(n,d) = 2\pi/T(n,d)$ is dimensionless (units of radians) and
is therefore not bounded below by any function of $\tP$.
For any $\varepsilon > 0$, there exists
\[
n_0 = 1 + \Bigl\lceil \log_{d+1}\!\Bigl(\frac{2\pi}{d\,\varepsilon}\Bigr)\Bigr\rceil
\]
such that $\Delta\theta(n,d) < \varepsilon$ for all $n \geq n_0$.
\end{theorem}

\begin{proof}
$\tP$ has dimension $[\tP] = \mathrm{s}$. The Heisenberg bound
$\Delta E \cdot \Delta t \geq \hbar/2$ constrains energy--time products;
both factors carry physical dimensions.
$\Delta\theta = 2\pi/T(n,d)$ is the ratio of two dimensionless quantities
and carries no time dimension: $[\Delta\theta] = \mathrm{rad} = 1$.
Neither the Heisenberg argument nor the gravitational collapse argument
applies to a dimensionless ratio.

Since $T(n,d) = d(d+1)^{n-1}$ is unbounded, $\Delta\theta \to 0$ as
$n \to \infty$. Setting $\Delta\theta < \varepsilon$ and solving gives $n_0$
as stated.
\end{proof}

\begin{remark}
The composition-inflation mechanism does not create shorter time intervals.
It creates more distinguishable trajectory structures within the same time
interval. The Planck limit bounds temporal duration;
composition-inflation operates on categorical multiplicity. These are
fundamentally different quantities.
\end{remark}

\subsection{The Planck depth}

\begin{definition}[Planck Depth]
\label{def:planckdepth}
The \emph{Planck depth} $n_P$ for a reference oscillator with period
$\tau_{\mathrm{osc}}$ in $d$-dimensional entropy space is the minimum
cycle count for which the categorical trajectory count equals or exceeds
the number of Planck-time intervals in one oscillator period:
\[
T(n_P, d) \geq \frac{\tau_{\mathrm{osc}}}{\tP}.
\]
\end{definition}

\begin{theorem}[Planck Depth Formula]
\label{thm:planckdepth}
\[
n_P = 1 + \Bigl\lceil \log_{d+1}\!\Bigl(\frac{\tau_{\mathrm{osc}}}{d \cdot \tP}\Bigr)\Bigr\rceil.
\]
\end{theorem}

\begin{proof}
Solve $d(d+1)^{n_P - 1} \geq \tau_{\mathrm{osc}}/\tP$ for $n_P$. \qed
\end{proof}

\begin{theorem}[Universality of Planck Depth]
\label{thm:universal}
For any physically realizable oscillator with frequency $\nu$ satisfying
$\tP^{-1} > \nu > 1\,\mathrm{Hz}$, the Planck depth at $d = 3$ lies in
$n_P \in [36, 73]$. The variation spans only 37 cycles despite frequencies
varying over 14 orders of magnitude, because $n_P$ depends logarithmically
on $\nu$.
\end{theorem}

\begin{proof}
At the lower frequency extreme ($\nu = 1\,\mathrm{Hz}$):
$\tau/\tP = 1.86 \times 10^{43}$,
$n_P = 1 + \lceil\log_4(1.86 \times 10^{43}/3)\rceil = 73$.
At the strontium optical clock ($\nu = 4.29 \times 10^{14}\,\mathrm{Hz}$):
$\tau/\tP = 4.32 \times 10^{28}$,
$n_P = 48$.
The logarithm compresses 14 orders of magnitude in frequency to
$73 - 48 = 25$ cycles.
\end{proof}

\begin{table}[h]
\centering
\caption{Planck depth for representative oscillators ($d = 3$).}
\label{tab:planck}
\begin{tabular}{lccc}
\toprule
Oscillator & $\nu$ (Hz) & $\tau/\tP$ & $n_P$ \\
\midrule
Sr optical clock & $4.29 \times 10^{14}$ & $4.3 \times 10^{28}$ & 48 \\
Cs-133 hyperfine & $9.19 \times 10^{9}$ & $2.0 \times 10^{33}$ & 56 \\
LHC bunch clock  & $4.01 \times 10^{7}$ & $4.6 \times 10^{35}$ & 60 \\
H maser          & $1.42 \times 10^{9}$ & $1.3 \times 10^{34}$ & 57 \\
CPU (3 GHz)      & $3.00 \times 10^{9}$ & $6.2 \times 10^{33}$ & 57 \\
\bottomrule
\end{tabular}
\end{table}

% =========================================================
\section{Categorical Angular Units and Constant Reduction}
\label{sec:constants}
% =========================================================

\subsection{The speed of light}

\begin{theorem}[Speed of Light as Angular Velocity]
\label{thm:c}
In categorical angular units, the speed of light is exactly
\begin{equation}
c = 2\pi\;\mathrm{rad/tick}.
\label{eq:c}
\end{equation}
\end{theorem}

\begin{proof}
In one oscillation cycle (one tick), a bounded system traverses the full
phase $\Delta\phi = 2\pi\,\mathrm{rad}$ of its reference oscillation.
The maximum rate at which categorical distinctions propagate between two
points separated by $\Delta x$ at the characteristic radius $r_0$
is $c = \Delta x/\tau_c = 2\pi r_0/\tau_{\mathrm{tick}}$ where
$\tau_{\mathrm{tick}} = 1\,\mathrm{tick}$. In angular units
$\theta = x/r_0$: $c = 2\pi\,\mathrm{rad/tick}$.

The partition propagation bound establishes that $c = \sup_n \ell_n \omega_n$
is finite and equals this value; the bound is attained at the massless
(zero-rest-frequency) limit.
\end{proof}

\subsection{Planck's constant}

\begin{theorem}[Planck's Constant as Energy per Radian]
\label{thm:hbar}
\begin{equation}
\hbar = E_{\mathrm{tick}}/(2\pi),
\label{eq:hbar}
\end{equation}
where $E_{\mathrm{tick}}$ is the energy carried by one oscillation tick.
\end{theorem}

\begin{proof}
From $E = \hbar\omega$ with $\omega = 2\pi/\tau_{\mathrm{tick}} = 2\pi\,\mathrm{tick}^{-1}$:
$E_{\mathrm{tick}} = \hbar \cdot 2\pi$. Solving: $\hbar = E_{\mathrm{tick}}/(2\pi)$.
$\hbar$ is the energy per radian of phase advance.
\end{proof}

\subsection{Boltzmann's constant}

\begin{theorem}[Boltzmann's Constant as Entropy per Composition Step]
\label{thm:kB}
\begin{equation}
k_B = E_{\mathrm{tick}}/\ln(d+1).
\label{eq:kB}
\end{equation}
\end{theorem}

\begin{proof}
By Proposition~\ref{prop:growth}(ii), each additional oscillation tick
multiplies the trajectory count by $(d+1)$, contributing $\ln(d+1)$ natural
units of categorical entropy. The thermal energy per degree of freedom is
$k_B T$, where $T$ is the tick rate. The energy per unit of composition
entropy is $E_{\mathrm{tick}}/\ln(d+1)$.
For $d = 3$: $k_B = E_{\mathrm{tick}}/\ln 4$.
\end{proof}

\subsection{Mass and rest energy}

\begin{corollary}[Mass and Rest Energy]
\label{cor:mass}
The inertial mass and rest energy of a localized oscillator with rest
frequency $\nu_0$ (ticks$^{-1}$) are
\begin{align}
m_0 &= \frac{E_{\mathrm{tick}} \cdot \nu_0}{4\pi^2},
\label{eq:mass}\\
E_0 &= m_0 c^2 = E_{\mathrm{tick}} \cdot \nu_0.
\label{eq:E0}
\end{align}
\end{corollary}

\begin{proof}
From $m_0 = \hbar\omega_0/c^2$ with $\hbar = E_{\mathrm{tick}}/(2\pi)$ and
$c = 2\pi\,\mathrm{rad/tick}$: $m_0 = E_{\mathrm{tick}}\omega_0/(2\pi \cdot 4\pi^2)
= E_{\mathrm{tick}}\nu_0/(4\pi^2)$ where $\nu_0 = \omega_0/(2\pi)$.
Then $E_0 = m_0 c^2 = E_{\mathrm{tick}}\nu_0/(4\pi^2)\cdot 4\pi^2 = E_{\mathrm{tick}}\nu_0$.
The identity $E_0 = \hbar\omega_0$ is recovered: rest energy is the fundamental
frequency-energy identity applied to the localization mode.
\end{proof}

\subsection{The Planck time and the irreducibility of $G$}

\begin{theorem}[Planck Time Decomposition]
\label{thm:tP}
\begin{equation}
\tP = \sqrt{\hbar G/c^5} = \frac{\sqrt{E_{\mathrm{tick}} \cdot G}}{(2\pi)^3}.
\label{eq:tP}
\end{equation}
The Planck time has two components: geometric factors $(2\pi)^3$ from
the angular structure of oscillation, and the product $\sqrt{E_{\mathrm{tick}} \cdot G}$
as the unique irreducible coupling.
\end{theorem}

\begin{proof}
Substitute $\hbar = E_{\mathrm{tick}}/(2\pi)$ and $c = 2\pi\,\mathrm{rad/tick}$:
$\tP = \sqrt{E_{\mathrm{tick}} G/(2\pi \cdot (2\pi)^5)} = \sqrt{E_{\mathrm{tick}} G}/(2\pi)^3$.
\end{proof}

\begin{theorem}[Irreducibility of $G$]
\label{thm:G}
The gravitational constant $G$ is the unique fundamental constant of physics
that does not reduce to a combination of tick counts, factors of $2\pi$,
and the energy unit $E_{\mathrm{tick}}$.
\end{theorem}

\begin{proof}
Table~\ref{tab:constants} shows the reduction of six of the seven
fundamental constants. The remaining constant $G$ appears only in
$\tP$ and the Schwarzschild radius; in both cases it appears as an
irreducible multiplicative factor that cannot be expressed in terms of
counting operations or geometry. It encodes the coupling between
categorical partition structure (which determines mass through
$m_0 = E_{\mathrm{tick}}\nu_0/(4\pi^2)$) and the curvature of space.
No counting argument reduces this coupling.
\end{proof}

\begin{table}[h]
\centering
\caption{Fundamental constants in categorical angular units.}
\label{tab:constants}
\begin{tabular}{lll}
\toprule
Constant & SI value & Categorical form \\
\midrule
$c$ & $3.00 \times 10^8\,\mathrm{m/s}$ & $2\pi\,\mathrm{rad/tick}$ \\
$\hbar$ & $1.055 \times 10^{-34}\,\mathrm{J\cdot s}$ & $E_{\mathrm{tick}}/(2\pi)$ \\
$k_B$ & $1.381 \times 10^{-23}\,\mathrm{J/K}$ & $E_{\mathrm{tick}}/\ln(d+1)$ \\
$m_0$ & $\hbar\omega_0/c^2$ & $E_{\mathrm{tick}}\nu_0/(4\pi^2)$ \\
$E_0$ & $m_0 c^2$ & $E_{\mathrm{tick}}\nu_0 = \hbar\omega_0$ \\
$\tP$ & $5.39 \times 10^{-44}\,\mathrm{s}$ & $\sqrt{E_{\mathrm{tick}}G}/(2\pi)^3$ \\
$G$ & $6.674 \times 10^{-11}\,\mathrm{m^3 kg^{-1} s^{-2}}$ & \textbf{irreducible} \\
\bottomrule
\end{tabular}
\end{table}

\begin{corollary}[Planck Depth from $G$]
\label{cor:planckdepthG}
Since $\tP = \sqrt{E_{\mathrm{tick}} G}/(2\pi)^3$ and $n_P$ depends on $\tP$
through Theorem~\ref{thm:planckdepth}, the Planck depth is computable to
precision $(d+1)^{-n}$ once $E_{\mathrm{tick}}$ and $G$ are known. Since $G$
is the unique irreducible constant, the Planck depth is the one quantity in
the entire categorical framework that is not derivable from geometry and
counting alone. It is the fingerprint of gravity in the composition-inflation
structure.
\end{corollary}

% =========================================================
\section{Discussion and Outlook}
\label{sec:discussion}
% =========================================================

\subsection{What has been established}

The paper has derived, from a single empirical axiom:

\begin{enumerate}[nosep,label=(\arabic*)]
\item Oscillatory necessity — bounded dynamics must oscillate
(Corollary~\ref{cor:oscillatory}).
\item The partition coordinate system $(n,\ell,m,s)$ with capacity $C(n) = 2n^2$
(Theorem~\ref{thm:capacity}).
\item The partition bijection $\Phi : \Zp \to \Parts$ and strict monotonicity
of the partition count (Theorems~\ref{thm:bijection}, \ref{thm:monotone}).
\item The composition-inflation formula $T(n,d) = d(d+1)^{n-1}$
(Theorem~\ref{thm:inflation}).
\item Angular resolution is dimensionless and Planck-unconstrained
(Theorem~\ref{thm:noplanck}).
\item The Planck depth formula and its universality across oscillators
(Theorems~\ref{thm:planckdepth}, \ref{thm:universal}).
\item All fundamental constants except $G$ reduce to $2\pi$, tick counts,
and $E_{\mathrm{tick}}$ (Theorems~\ref{thm:c}--\ref{thm:G} and Table~\ref{tab:constants}).
\item $G$ is the unique irreducible constant (Theorem~\ref{thm:G}).
\end{enumerate}

\subsection{Applications developed in companion papers}

The partition algebra derived here is the mathematical foundation for a
sequence of companion papers. The following applications are developed
there, citing only results proved above:

\textbf{LHC trigger design.} A trigger system with $d$ independent detector
subsystems observing $n$ bunch crossings can discriminate $T(n,d)$ distinct
event categories. The Planck depth $n_P = 60$ for the LHC bunch clock
($\nu_B = 40\,\mathrm{MHz}$, $d = 3$) defines a natural design depth;
the current L1 trigger at $n = 100$ operates at $n/n_P = 1.67$,
placing it in the super-Planck categorical resolution regime
\citep{companion:fundamentaltiming,companion:planckdepth}.

\textbf{Mass spectrometry.} The partition Lagrangian
$\mathcal{L}_\mathcal{M} = \frac{1}{2}\mu|\dot{\mathbf{x}}|^2 +
\mu\dot{\mathbf{x}}\cdot\mathbf{A}_\mathcal{M} - \mathcal{M}(\mathbf{x},t)$
derives all four mass analyzer equations (TOF, quadrupole, Orbitrap, FT-ICR)
from the same Euler-Lagrange equations, with the Lorentz force as a
consequence rather than a postulate \citep{companion:ion}.

\textbf{Gravitational constant.} Three independent routes to $G$ via the
partition algebra — oscillation-ratio, category fixed-point, and
partition-density ratio — converge on the same value with
$(d+1)^{-n}$ precision. At $n = 8$ all three agree with CODATA;
at $n = 27$ they reach double-precision machine epsilon
\citep{companion:G}.

\textbf{Cosmology.} A partition-density-dependent $G$ reproduces the MOND
acceleration scale $a_0 = cH_0/(2\pi)$, predicts dark energy equation of state
$w_{\mathrm{eff}} = -1 + 1/(d+1) = -0.75$, and secular drift
$\dot G/G = -3H_0/(d+1)$ \citep{companion:cosmology}.

\subsection{Falsification conditions}

The framework is falsified by any of the following:
\begin{itemize}[nosep]
\item A persistent physical system with infinite Liouville measure.
\item A bounded oscillatory system in which the partition count $M(t)$
decreases (i.e., a partition decrement is physically realized).
\item A physical observable that cannot be expressed as a function of
$(n,\ell,m,s,E_{\mathrm{tick}},G)$.
\end{itemize}
None of these conditions has been observed.

% =========================================================
\section*{Acknowledgements}
This work was supported by the AIMe Registry for Artificial Intelligence.
% =========================================================

\bibliographystyle{plainnat}
\bibliography{references}

\end{document}
